\chapter{Assembling the transformer blocks}
\label{ch:transformer_block}

In this chapter,\keyterm{Transformer block (transformer block)}is created by combining the attention from Chapter 5 and the embedding from Chapter 4. Attention alone is not enough; only a combination of FFN (Feed-Forward Network), residual connection, and layer normalization can learn meaningful expressions.

\section{Structure of transformer block}

One transformer block consists of four components:

\begin{enumerate}
\item\textbf{Multi-Head Attention}: Information exchange between tokens
\item\textbf{Feed-Forward Network (FFN)}: Nonlinear transformation by position
\item\textbf{Residual Connection}: Deep network training by adding input to output
\item\textbf{Layer Normalization}: Stabilization of learning
\end{enumerate}

How these four things are combined is the key to the transformer block. Volume 1 uses the Pre-LN structure (GPT-2 style).

\begin{figure}[htbp]
\centering
\begin{tikzpicture}[
  every node/.style={font=\small\sffamily},
  box/.style={draw=inkblue, fill=inkblue!8, rounded corners=2pt, minimum width=3cm, minimum height=0.7cm},
  norm/.style={draw=inkgreen, fill=inkgreen!10, rounded corners=2pt, minimum width=2.5cm, minimum height=0.6cm},
  sum/.style={draw=inkred, fill=inkred!10, circle, minimum size=0.6cm}
]
  \node[box] (in) at (0, 0) {input $x$};
  \node[norm] (ln1) at (0, -1) {LayerNorm};
  \node[box] (attn) at (0, -2) {Multi-Head Attention};
  \node[sum] (add1) at (0, -3) {$+$};
  \node[norm] (ln2) at (0, -4) {LayerNorm};
  \node[box] (ffn) at (0, -5) {FFN (GELU)};
  \node[sum] (add2) at (0, -6) {$+$};
  \node[box] (out) at (0, -7) {output};

  % residual bypass line
  \draw[inkblue, rounded corners=4pt] (-2, 0) -- (-2, -3);
  \draw[-{Stealth[length=4pt]}, inkblue] (-2, -3) -- (add1);
  \draw[inkblue, rounded corners=4pt] (2, -3) -- (2, -6);
  \draw[-{Stealth[length=4pt]}, inkblue] (2, -6) -- (add2);

  % main flow
  \draw[-{Stealth[length=4pt]}, inkblue] (in) -- (ln1);
  \draw[-{Stealth[length=4pt]}, inkblue] (ln1) -- (attn);
  \draw[-{Stealth[length=4pt]}, inkblue] (attn) -- (add1);
  \draw[-{Stealth[length=4pt]}, inkblue] (add1) -- (ln2);
  \draw[-{Stealth[length=4pt]}, inkblue] (ln2) -- (ffn);
  \draw[-{Stealth[length=4pt]}, inkblue] (ffn) -- (add2);
  \draw[-{Stealth[length=4pt]}, inkblue] (add2) -- (out);

  % label
  \node[inkgray, font=\scriptsize, anchor=west] at (-3.5, -1.5) {residual};
  \node[inkgray, font=\scriptsize, anchor=west] at (2.5, -4.5) {residual};
\end{tikzpicture}
\caption{Pre-LN Transformer block structure. Attention and FFN Connect the residuals to each and LayerNormThere is.}
\label{fig:transformer-block}
\end{figure}

\section{Feed-Forward Network (FFN)}

Attention performs ``horizontal'' information exchange between tokens. FFN, on the other hand, performs a ``vertical'' non-linear transformation at each token position. The structure is simple:

\[
\text{FFN}(x) = \text{Linear}\_2(\text{GELU}(\text{Linear}\_1(x)))
\]

$\text{Linear}\_1$expands to$d\_{model} \to d\_{ff}$, and$\text{Linear}\_2$contracts back to$d\_{ff} \to d\_{model}$.$d\_{ff}$is usually$4 d\_{model}$, and this ``expand-contract'' pattern gives non-linear expressiveness.

\subsection{intuition: why expand-Is it shrinking??}

The reason for expanding$d\_{model} = 512$to$d\_{ff} = 2048$and then contracting it back to 512 is to ``operate in a larger expression space.'' In higher-dimensional spaces, more complex decision boundaries can be drawn. Combined with the non-linearity of the activation function (GELU), FFN learns more than just a linear transformation.

In fact, research results show that FFN acts as a ``key-value memory'' (Geva et al. 2021).  The interpretation is that the row vector of$\text{Linear}\_1$acts as the ``key'' and the column vector of$\text{Linear}\_2$acts as the ``value'', so FFN is actually an information retrieval mechanism.

\begin{lstlisting}[language=Python]
import torch.nn.functional as F

class FeedForward(nn.Module):
    def __init__(self, d_model: int, d_ff: int, dropout: float = 0.1):
        super().__init__()
        self.w1 = nn.Linear(d_model, d_ff)
        self.w2 = nn.Linear(d_ff, d_model)
        self.dropout = nn.Dropout(dropout)
        self._init_weights()

    def _init_weights(self) -> None:
        nn.init.normal_(self.w1.weight, mean=0.0, std=0.02)
        nn.init.normal_(self.w2.weight, mean=0.0, std=0.02)
        if self.w1.bias is not None:
            nn.init.zeros_(self.w1.bias)
        if self.w2.bias is not None:
            nn.init.zeros_(self.w2.bias)

    def forward(self, x: torch.Tensor) -> torch.Tensor:
        """x: [batch, seq, d_model] -> [batch, seq, d_model]"""
        return self.dropout(self.w2(F.gelu(self.w1(x))))
\end{lstlisting}

\section{GELU vs ReLU}

The original Transformer used ReLU, but since GPT-2,\keyterm{GELU(Gaussian Error Linear Unit)}is the standard. GELU uses a Gaussian cumulative distribution for smooth transitions instead of ReLU's ``harsh'' 0/1 threshold.

\begin{formula}[GELU]
\[
\text{GELU}(x) = x \cdot \Phi(x) \approx 0.5 x \left(1 + \tanh\left[\sqrt{2/\pi}(x + 0.044715 x^3)\right]\right)
\]
$\Phi(x)$is the cumulative distribution function of the standard normal distribution.$x < 0$also allows small negative values, so the gradient flow is better.
\end{formula}

\begin{figure}[htbp]
\centering
\begin{tikzpicture}[
  scale=0.8,
  every node/.style={font=\small\sffamily}
]
  % axis
  \draw[inkblue, line width=0.5pt, ->] (-3, 0) -- (3, 0) node[right] {$x$};
  \draw[inkblue, line width=0.5pt, ->] (0, -0.5) -- (0, 2.5) node[above] {activation};
  % ReLU (Bend)
  \draw[inkred, line width=1pt] (-3, 0) -- (0, 0) -- (2.5, 2.3);
  \node[inkred, font=\scriptsize] at (2, 2.3) {ReLU};
  % GELU (soft curve)
  \draw[inkblue, line width=1pt, domain=-3:2.5, samples=50] plot (\x, {0.5*\x*(1+tanh(0.7978*(\x+0.044715*\x*\x*\x)))});
  \node[inkblue, font=\scriptsize] at (1.5, 1.5) {GELU};
\end{tikzpicture}
\caption{ReLU vs GELU. ReLUIs   --  is exactly 0, but, GELUIs 작은 음수값을 허용하여 부드럽게 전환.}
\label{fig:gelu-relu}
\end{figure}

\subsection{why GELUIs it better??}

ReLU becomes completely 0 at$x < 0$, so the gradient becomes 0 (``dead ReLU'' problem). GELU allows small negative values ​​so the gradient always flows. Additionally, the transition point is smooth and differentiable, making optimization stable.

In fact, most modern models, including GPT-2, BERT, and RoBERTa, use GELU. Volume 2 covers a more advanced activation function called SwiGLU.

\section{Connect residuals (Residual Connection)}

\keyterm{Concatenate residuals (residual connection)}is an operation that adds input directly to output:$y = x + f(x)$. This simple idea made deep network learning possible. He et al. (2016) started with ResNet, which enabled learning over 100 layers in image classification.

\subsection{intuition: Why residuals help?}

If there are no residuals, the network must learn$f(x) = y$. The deeper you go, the more difficult it is to learn this mapping. If there are residuals, you only need to learn$f(x) = y - x$, that is, ``residuals''. If the input$x$is already close to$y$,$f$only needs a small correction.

Additionally, when backpropagating, the gradient flows directly through a ``shortcut''.  The derivative of$y = x + f(x)$is$\frac{dy}{dx} = 1 + f'(x)$.  Even if it is$f'(x) = 0$, the gradient becomes 1 and does not disappear.

The transformer block has two residual connections: one around Attention and one around FFN. GPT has 12$\sim$96 layers, so learning is impossible without residuals.

\begin{warnbox}[100-layer network without residual]
Experimentally, if you subtract the residuals and train a 100-layer network, the loss does not decrease. This is because the gradient disappears during backpropagation and the front layer is not learned. Residual connections are simple but essential elements of deep networks.
\end{warnbox}

\section{Pre-LN vs Post-LN}

There are two variations depending on the location of the LayerNorm.

\subsection{Post-LN (Vaswani text)}

The original Transformer placed the LayerNorm at Attention/FFN\emph{back}:

\[
x = \text{LN}(x + \text{Attn}(x))
\]

However, this structure is unstable in learning in deep networks. It is very sensitive to initialization and prone to divergence without warmup.

\subsection{Pre-LN (GPT-2 styles)}

Place LayerNorm in Attention/FFN\emph{front}:

\[
x = x + \text{Attn}(\text{LN}(x))
\]

This structure has no LayerNorm in the residual path, so the gradient flows better. It is learned stably even in deep networks without warmup, making it the standard for modern LLM.

\subsection{why Pre-LNIs this more stable??}

After going through the residual path$x \to x + \text{Attn}(x) \to \text{LN}(\cdot)$in Post-LN, the gradient must pass through LayerNorm. LayerNorm is an operation that divides by variance, so it scales the gradient, but if this scaling is accumulated in a deep network, the gradient may disappear.

In Pre-LN,$x$flows directly from the residual path$x \to x + \text{Attn}(\text{LN}(x))$without going through LayerNorm. Only the attention output is added after going through LayerNorm. Learning is possible even in deep networks as the gradient flows directly through the identity path.

\begin{lstlisting}[language=Python]
class TransformerBlock(nn.Module):
    def __init__(self, config: ModelConfig, layer_norm_style: str = "pre"):
        super().__init__()
        self.layer_norm_style = layer_norm_style
        self.ln1 = nn.LayerNorm(config.d_model, eps=config.layer_norm_eps)
        self.ln2 = nn.LayerNorm(config.d_model, eps=config.layer_norm_eps)
        self.attn = MultiHeadAttention(config)
        self.ffn = FeedForward(config.d_model, config.d_ff, config.dropout)

    def forward(self, x, mask=None):
        if self.layer_norm_style == "pre":
            # Pre-LN: normalization -> Attention -> residual
            x = x + self.attn(self.ln1(x), mask=mask)
            x = x + self.ffn(self.ln2(x))
        else:
            # Post-LN: Attention -> residual -> normalization
            x = self.ln1(x + self.attn(x, mask=mask))
            x = self.ln2(x + self.ffn(x))
        return x
\end{lstlisting}

\begin{warnbox}[Reason for using Pre-LN]
Post-LN can disappear in deep networks because the initial gradient flows only through the LayerNorm and not the residual path. In Pre-LN, there is nothing in the residual path, so the gradient flows identity, making learning possible even with more than 100 layers. However, Pre-LN requires an additional LayerNorm at the end (implemented as\code{ln\_f}in the next chapter).
\end{warnbox}

\section{Residual scaling}

GPT-2 initializes the weights that directly contribute to the residual stream ($W\_O$in attention,$W\_2$in FFN) to be smaller. This is to prevent the dispersion of residual output from increasing as the layer becomes deeper.

\begin{lstlisting}[language=Python]
def _init_residuals(self, module: nn.Module) -> None:
    """Initialize projection weights that contribute directly to the residual stream to be small.."""
    if isinstance(module, nn.Linear):
        if module.weight.shape[0] == self.config.d_model:
            # Output to residual: Initialize to a smaller standard deviation
            nn.init.normal_(module.weight, mean=0.0,
                            std=0.02 / math.sqrt(2 * self.config.n_layers))
\end{lstlisting}

$1/\sqrt{2N}$scaling corrects the fact that the variance of the residual output increases in proportion to the number of layers when there are$N$layers.  The larger$N$, the smaller initialization is needed.

\subsection{why   --  impression?}

Each transformer block has two outputs (attention, FFN) that contribute to the residual.  If there are$N$blocks, a total of$2N$outputs are added to the residual. If the variance of each output is$\sigma^2$, the variance of the sum is$2N \sigma^2$. For this to become 1, it must be$\sigma = 1/\sqrt{2N}$.

\section{block test}

\begin{lstlisting}[language=Python]
def test_block_shape(small_config):
    block = TransformerBlock(small_config)
    x = torch.randn(2, 8, small_config.d_model)
    out = block(x)
    assert out.shape == (2, 8, small_config.d_model)

def test_block_residual_connection(small_config):
    """If there is a residual connection, the difference between input and output must be finite."""
    block = TransformerBlock(small_config)
    block.eval()
    x = torch.randn(2, 8, small_config.d_model)
    out = block(x)
    diff = (out - x).abs().mean()
    assert diff.item() < 5.0  # Normalization problems if residuals are too large

def test_feedforward_shape(small_config):
    ffn = FeedForward(small_config.d_model, small_config.d_ff)
    x = torch.randn(2, 8, small_config.d_model)
    out = ffn(x)
    assert out.shape == x.shape
\end{lstlisting}

\section{Computation amount of transformer block}

The FLOPs of one block are approximately:
\begin{itemize}
\item QKV projection:$3 \cdot L \cdot d^2$
\item Attention Score:$L^2 \cdot d$
\item Attention output projection:$L \cdot d^2$
  \item FFN: $2 \cdot L \cdot d \cdot d\_{ff} = 8 L d^2$ (since $d\_{ff} = 4d$)
\end{itemize}

The total is$O(L \cdot d^2)$, and$O(L^2 d)$of attention can be ignored when it is$L < d$. In actual GPT-3 ($d = 12288$,$L = 2048$), FFN is more dominant because it is$L \cdot d^2 = 3 \times 10^{11}$and$L^2 \cdot d = 5 \times 10^{10}$.

\section{summation}

\begin{itemize}
\item Transformer block = Attention + FFN + Residual connection + LayerNorm.
\item FFN is a non-linear transformation by position (Linear$\to$GELU$\to$Linear).
\item GELU improves learning stability with smoother transitions than ReLU.
\item Residual connection is the core of deep network learning.
\item Pre-LN has better learning stability than Post-LN, a modern LLM standard.
\item Residual scaling$1/\sqrt{2N}$to prevent distributed explosions in deep networks.
\end{itemize}

\section{practice problems}

\begin{enumerate}
\item What happens if$d\_{ff}$in FFN is changed to$d\_{model}$instead of$4d\_{model}$?
\item What phenomenon is observed when residual connections are removed and a 12-layer model is trained?
\item Why is an extra LayerNorm needed after the last block in Pre-LN?
\item What is the difference if I use Sigmoid Linear Unit (SiLU,$x \cdot \sigma(x)$) instead of GELU?
\item If we initialize a 100-layer model without residual scaling, what happens at the beginning of training?
\end{enumerate}

In the next chapter, we build a complete GPT model by stacking these blocks$N$layers and build a learning loop.
